\input{preamble}

\begin{document}

\title{Springer theory}
\maketitle

\phantomsection
\label{section-phantom}

\section{Springer theory}
\label{section-springer-theory}

\noindent
Simple Lie algebra $\mathfrak{g}$. Nilpotent
cone $\mathcal{N} \subset
\mathfrak{g}$,
$\mathcal{N}=\{ x \in \mathfrak{g}:
\text{ad}(x)\text{ nilpotent}\}$
[we've already seen this object].

$G$ (connected complex) Lie group with Lie
algebra $\mathfrak{g}$. There is the adjoint
action of $G$ on $\mathfrak{g}$. (If $G
\subset \text{Mat}_{n\times n}(\mathbb{C})$,
and $\mathfrak{g}= \text{Lie} G\subset
\text{Mat}_{n \times n}(\mathbb{C})$,
$G$ acts on $\mathfrak{g}$ by conjugation,
$\text{ad}(g)\cdot X:=g X g^{-1}$.)

$$
\mathfrak{g}=\coprod (\text{infinitely many
adjoint orbits}),
$$
but
\begin{equation}
\label{s-label-001}
\mathcal{N}=\coprod(\text{finitely many
adjoint orbits}).
\end{equation}

\noindent
Equation
\ref{s-label-001}
is proved via a sequence of results:
\begin{enumerate}
\item Jacobian-Morozov theorem
(each nilpotent $e \in \mathfrak{g}$ is
contained in an $\mathfrak{sl}_2$-triple
$\{e,h,f\}$),

\item ``Conjugacy theorems of (a) Kostant
and (b) Malcer''.

\item Classification of the possible
orbits of $h$ in the $\mathfrak{sl}_2$-triple
$\xymatrix{\ar@{~>}[r]&}$ wieghed Dynkin
diagrams.
\end{enumerate}

\medskip\noindent
Springer correspondence:
$$
\text{Irreps of $W$}
\leftrightarrow
\text{Nilpotent orbits}.
$$
\begin{example}
\label{example-sln-nilpotent}
$\mathfrak{g}=\mathfrak{sl}_n$, then nilp $e
\in \mathfrak{sl}_n$. Put $e$ in Jordan
canonical form
$$
\left[\substack{\text{Matrix decomposed in
Jordan blocks,}\\
\text{each block with all zero but 1's over
the diagonal}\\
\text{- because it's nilpotent it looks like
that}}\right]
$$
That is, a partition of $n$. $[n]$ correspond
to orbit called principal nilpotent orbit (or
``regular nilpotent orbit''). In particular
$[1^n]$ corresponds to $e=0$.

So \# of nilpotent orbits = $p(n) $ =
\# integer partitions of $n$.

For $\mathfrak{sl}_n$, Weyl group is $W
\simeq S_n$. A conjugacy class in $S_n$ is
given by ``cycle shape'' determined by
lengths of cycles. [Example.]

If $G$ finite group, we know that
\# conjugacy classes = \# irreps.

For $S_n$, this also = $p(n)$. [In general
it's not clear how to assign conjugacy
classes to irreps; $S_n$ is an exception]
\end{example}

\noindent
For $\mathfrak{g}=E_8$, $\# W \simeq 700
\times 10^6$,
\# conjugacy classes in  $W=112$.
\# nilpotent orbits in $\mathcal{N}(E_8)=70$.

\medskip\noindent
Key objects. Start with $\mathfrak{g}$. A
Borel subalgebra $\mathfrak{b}\subset
\mathfrak{g}$
is a maximal (wrt inclusion) solvable
subalgebra. Denote by $\mathcal{B}$ (or
$\mathbb{B}$) the set of all Borels in
$\mathfrak{g}$.

\noindent
Facts. By the adjoint action of $G$, all
Borels are conjugate. ($G$ acts transitively
in $\mathbb{B}$.)

This $\implies \mathbb{B}
\overset{\text{sets}}{\simeq}G/
Stab(\mathfrak{b}_0)
\subset G$
the Borel subgroup corresp.

For $\mathfrak{sl}_n$, can take
$$
\mathfrak{b}+0=\left\{
\begin{pmatrix} \ast & \ldots{} & \\
& \ast & \ast \\
& & \ddots \\
0 & & \ldots{} & \ast\end{pmatrix} \right\}
$$
(upper triangular matrices).

And $B_0=$ same but in $\text{SL}_n$
invertible, det=1, upper triangular

\noindent
Aside. If $\mathfrak{g}$ has rank $\ell$,
root system  $\Delta$,
$\mathfrak{g}=\mathfrak{n}_-\oplus
\mathfrak{h}\oplus \mathfrak{n}_+$
and $\mathfrak{b}=\mathfrak{h} \oplus
\mathfrak{n}_+$
is an example of a Borel algebra. Dimension
is $m=\ell+\frac{1}{2}|\Delta|$.

\medskip\noindent
Back to main point. Since all Borels are
conjugate, they all have the same dimension.
$\mathbb{B}\subset
\text{Gr}(\mathfrak{g},m)$,
the Grassmannian. $\mathbb{B}\subset
\text{Gr}(\mathfrak{g},m)$ is an algebraic
subvariety [which may not be smooth]. Since
$\mathbb{B}$ has a transitive $G$-action, a
local homeo with $\mathbb{C}^m$ near any
point $\implies$ local homeos everywhere, so
$\mathbb{B}$ is smooth. This $\mathbb{B}$ is
called the flag variety $\mathbb{B} \simeq
G/B$.

Over $\text{Gr}(V,m)$ in general one has the
{\it tautological bundle}, i.e. over the
$m$-dimensional vector space $U \subset V$,
over the point $[U] \in \text{Gr}(V,m)$ one
`puts' $U$ itself.

\medskip\noindent
We do the same for $\mathbb{B}$. Define
$\tilde{\mathfrak{g}}=
\{(\mathfrak{b},x):\mathfrak{b} \in
\mathbb{B}, x \in \mathfrak{b}\}$.
Then $(\mathfrak{b},x) \mapsto  \mathfrak{b}$
defines $\tilde{\mathfrak{g}}\to \mathbb{B}$
and $\tilde{\mathfrak{g}}$ is a sort of
tautological vector bundle of $B$.

Can also map
$$
\begin{aligned}\tilde{\mathfrak{g}}
&\longrightarrow \mathfrak{g} \\
(\mathfrak{b},x) &\longmapsto x
\end{aligned}
$$
$$
\xymatrix{
&\tilde{\mathfrak{g}}\ar[dl]_{\mu}\ar[dr]
^{\pi}\\
\mathfrak{g}&&\mathbb{B}.
}
$$

\medskip\noindent
\{It looks like the projection $\mu$
is more interesting than $\pi$: fix an
element of the Lie algebra and look at all
the Borels that contain this element; indeed,
0 is contained in every Borel and thus we get
the ``zero section'' on top of 0.\}

\begin{definition}
\label{definition-springer-fibre}
Let $x \in \mathfrak{g}$, the subset
$\mu^{-1}(x) \subset
\tilde{\mathfrak{g}}$
is called a {\it Springer fibre}.
\end{definition}

\begin{example}
\label{example-flag-variety}
$\mu^{-1}(0)=\mathbb{B}$ [since all borels
contain zero, this is kind of a zero section
of the bundle, i.e. a copy of $\mathbb{B}$],
the flag variety.
\end{example}
Other extreme: Let $x \in \mathfrak{g}$ be a
regular semisimple element. Then
$$
|\mu^{-1}(x)|=|W|,
$$
finite set in bijection w/ Weyl group.

$$
\left[
\substack{\text{Picture of plane containing
the nilpotent cone $\mathcal{N}$} \\
\text{with $\mathbb{B}$ sitting over
origin,}\\
\text{a fiber of 4 points on top a point
labelled ss reg}\\
\text{and other varieties on top of other
points in the base plane.}}\right]
$$
$h \in \mathfrak{g}$ regular semisimple, then
$h$ sits in a unique Cartan $\mathfrak{h}$,
then $\mathfrak{h}$ sits inside $|W|$
disjoint Borels.

\begin{remark}
\label{remark-pincipal-orbit}
For $x \in \mathbb{O}_{\text{prin}}\subset
\mathcal{N}$
(principal orbit), $\mu^{-1}(x)=(\text{one
point})$. But for $x$ in smaller orbits,
$\mu^{-1}(x)$ becomes a higher dimensional
variety. In fact
\begin{equation}
\label{equation-important}
\dim\mu^{-1}(x)=\frac{1}{2}(\dim
\mathcal{N}-\dim \mathbb{O})
\end{equation}
for $x \in \mathbb{O}$.
\end{remark}

\medskip\noindent
Consider $\mu|_{\mu^{-1}(\mathcal{N})}$ i.e.
$$
\tilde{\mathcal{N}}=\{(\mathfrak{b},x):
\mathfrak{b}
\in \mathbb{B},
x \in \mathfrak{b}\text{ nilpotent}\}.
$$
[Now we consider a diagram as before but
restricted (and abuse notation denoting the
maps with the same letters:]
$$
\xymatrix{
&
\tilde{\mathcal{N}}\ar[dr]^{\pi}\ar[dl]_{\mu}
\\
\mathcal{N} & & \mathbb{B}.
}
$$
If $\mathfrak{g}=\mathfrak{n}_-+\mathfrak{h}+
\mathfrak{n}_-$,
$x \in \mathfrak{b}$ nilpotent $\iff x \in
\mathfrak{n}
=[\mathfrak{b},\mathfrak{b}]$,
$\mathfrak{b}=\mathfrak{h}+\mathfrak{n}_-$,
$[\mathfrak{b},\mathfrak{b}]=\mathfrak{n}_-$.

Notice that
$\tilde{\mathcal{N}}\xrightarrow{\pi}
\mathbb{B}$
is also a vector bundle [fibres are vector
spaces of the same dimension]. Now fibre over
$\mathfrak{b}$ is $\mathfrak{n}\subset
\mathfrak{b}$.

Claim.
$$
\xymatrix{
\tilde{\mathcal{N}}\ar[d]\\\mathbb{B}
}
$$
is isomorphic to the cotangent bundle.
$\tilde{\mathcal{N}}=T^*_{\mathbb{B}}$
(bundle of 1-forms).
$$
\dim \mathbb{B}
=\dim G - \dim B
=\dim \mathfrak{g}- \dim \mathfrak{b}.
$$
[It's not hard to see this bundle is indeed
the cotangent bundle, for instance we check
that dimension must be
$$
\begin{aligned}
\dim \mathfrak{g}+\dim \mathfrak{b}
&=\dim \mathfrak{n}_-\\
&=\dim \mathfrak{n}_+\\
&=\text{rk}(\tilde{\mathcal{N}})
\end{aligned}
$$
since $\mathfrak{g}
=\mathfrak{n}_-+\underbrace{\mathfrak{h}+
\mathfrak{n}_+}_{\mathfrak{b}}$.]

[Kaledon's theory - when we pushforward a
constant sheaf on a resolution of varieties
where the upstairs variety has symplectic
structure. See perverse sheaves, Equation
\ref{equation-important}.]

\medskip\noindent
Let $h$ be regular semisimple, $h \in
\mathfrak{h}$ (we know that $\mathfrak{h}$ is
unique). How many $\mathfrak{b}$ Borels
contain $\mathfrak{h}$? Recall $N_G(T)/T
\simeq W$ [$T$ torus].

\medskip\noindent
Let $\mathfrak{b}_1,\mathfrak{b}_2$ be two
Borels, $h \in
\mathfrak{b}_1,\mathfrak{b}_2$. There exists
$g$ such that
$g\mathfrak{b}g^{-1}=\mathfrak{b}$.

\begin{remark}
\label{remark-flag}
$G$ acts on
$\mathbb{B}\xrightarrow{\simeq}G/B$
$$
\xymatrix{
\mathfrak{b}\ar@{|->}[d]\ar[r]&
g_0B\ar@{|->}[d]\\
g\mathfrak{b}g^{-1}\ar[r]&g g_0B.
}
$$
Consider the restricted action of $B$ on
$G/B$ and its orbits. $B\setminus G/B$
``double cosets''.
\end{remark}

\noindent
Fact. There exist finitely many double
cosets, in fact $|W|$. Indeed
$$
\xymatrix@R=0.1em{
W\ar@{-}[r]
&NG(T)/T\ar@{-}[r]
&  B\setminus G / B \ar@{-}[r]
& \mathbb{B}\times \mathbb{B}/\Delta G\\
&g \ar@{|->}[r]& g\\
& & BgB \ar@{|->}[r]
&(gB,eB)\text{ or something}
}
$$
Recall that $T \subset B \subset G$.
[Concerning the rightmost group: the action
if $g(\mathfrak{b}_1,\mathfrak{b}_2)
=(g\mathfrak{b}_1g_1,g\mathfrak{b}_2g^{-1})
$.]

[$W$ can be simply thought of as the group
generated by the reflections determined by
the roots; $N_G(T)/T$ is similar, recall that
$T$ is the\ldots{}]

\medskip\noindent
[We shall see one last variety, the Steinberg
variety.]

Recall
$$
\xymatrix{
\tilde{\mathcal{N}}=
\{(\mathfrak{b},x):x \in \mathcal{N},
\mathfrak{b} \ni x\}
\ar[d]^{\mu}\\
\mathcal{N}.
}
$$
Define [a fibered product]
$$
\xymatrix{
Z=\tilde{\mathcal{N}}\times_{\mathcal{N}}
\tilde{\mathcal{N}}
=- \{(\mathfrak{b}_1,\mathfrak{b}_2,x):
x \in \mathcal{N},
\mathfrak{b}_1,\mathfrak{b}_2 \ni x\}
\ar[d]^{\mu^{(2)}}\\
\mathcal{N}.
}
$$
If $\mathbb{O}\subset \mathcal{N}$ is a
nilpotent orbit, $x \in \mathbb{
\mathcal{O}}$,
$\dim \mu^{-1}(x)=\frac{1}{2}(\dim
\mathcal{N} - \dim \mathbb{ O})$.

$Z$ is called the {\it Steinberg variety}.

$$
\xymatrix{
\mu^{-1}(\mathbb{O})\ar[d]
^{\substack{\text{fibration}
\\ \mu}}\\
\mathbb{O}
}
$$

\medskip\noindent
Let $Y_W \subset \mathbb{B} \times
\mathbb{B}$,
the $\Delta G$-orbit correspond to $w \in W$.
Define $Z_W = T^*_{Y_W}(\mathbb{B} \times
\mathbb{B})$
the  ``conormal bundle'' of $Y_n \subset
\mathbb{B} \times \mathbb{B}$
by, $C_p=\{ \varphi \in T^*_pM:
\varphi(T_p N)=0\}$

[We have two ways of partitioning $Z$, one is
by in which orbit an elements lies, i.e.
using
$$
\xymatrix{
Z\ar[d]^{\mu^{(2)}}\\
\mathbb{O},
}
$$
(observe that dimension = $\dim
\mathbb{O}+2\cdot \frac{1}{2}
(\dim \mathcal{N}-\dim \mathbb{ O}=\dim
\mathcal{N}$.


and another using parametrized by the Weyl
group as in Equation
\ref{equation-partition-w}

\medskip\noindent
Claim.
$$
Z \cong \coprod_{w \in W}Z_w
$$


\section{Summary of lecture 1}
\label{section-summary-1}

\begin{itemize}
\item Lie group $G$ with lie algebra
$\mathfrak{g}$. The set of Borel subalgebras
of $\mathfrak{g}$ is a smooth algebraic
variety $\mathbb{B}$.

\item Define $\tilde{\mathfrak{g}}=
\{(\mathfrak{b},x):\mathfrak{b} \in
\mathbb{B}, x \in \mathfrak{b}\}$.
Then $(\mathfrak{b},x) \mapsto  \mathfrak{b}$
defines $\tilde{\mathfrak{g}}\to \mathbb{B}$
and $\tilde{\mathfrak{g}}$ is a sort of
tautological vector bundle of $B$.

Can also map
$$
\begin{aligned}\tilde{\mathfrak{g}}
&\longrightarrow \mathfrak{g} \\
(\mathfrak{b},x) &\longmapsto x
\end{aligned}
$$
$$
\xymatrix{
&\tilde{\mathfrak{g}}\ar[dl]_{\mu}\ar[dr]
^{\pi}\\
\mathfrak{g}&&\mathbb{B}.
}
$$

\medskip\noindent
\{It looks like the projection $\mu$
is more interesting than $\pi$: fix an
element of the Lie algebra and look at all
the Borels that contain this element; indeed,
0 is contained in every Borel and thus we get
the ``zero section'' on top of 0.\}

\item Restrict those maps to
$$
\tilde{\mathcal{N}}=\{(\mathfrak{b},x):
\mathfrak{b}
\in \mathbb{B},
x \in \mathfrak{b}\text{ nilpotent}\}.
$$

\item $\tilde{\mathcal{N}}$ is the cotangent
bundle of $\mathbb{B}$!
\end{itemize}




\end{document}
